RQ3 asks how robust the RQ1/RQ2 findings are to reasonable alternative
analytical choices, and whether they hold across compensation groups. We
address this with two independent exercises: a specification-vibration
analysis across many plausible model variants, and a permutation-based
multi-group analysis (MGA) comparing compensated and uncompensated
residents directly. As the StaBLab consulting-methodology notes that guided
this project put it, the goal of a robustness exercise is not to search for
significant results, but to examine how sensitive the main findings are to
different reasonable analytical choices -- trying many specifications and
reporting only the one that happens to be significant would be a form of
selective reporting, not robustness testing.

\subsection{Specification vibration}
\label{robustness:vibration}

We first compared six theoretically motivated model specifications, each
changing exactly one analytical choice at a time (e.g.\ three-item vs.\
four-item tolerance, alternative cost definitions, alternative control sets).
We then went further and estimated every possible combination of 13 binary
adjustment choices -- $2^{13} = 8{,}192$ models in total -- following the
vibration-of-effects framework of \citet{Patel2015}, who found that
observational associations can flip sign (a "Janus effect") under different,
individually defensible specifications; in their reanalysis of 417
variables, 31\% showed such a sign flip. Across all 8{,}192 model variants
estimated here, the compensation participation -- tolerance association
was negative in every single model (0 Janus effects), and reached
conventional significance ($p < .05$) in 99.4\% of them. The median
standardized coefficient across all 8{,}192 models was $\beta = -0.09$,
consistent with the range obtained across the six theoretically motivated
specifications ($\beta = -0.128$ to $-0.143$) and with the primary model
reported in Section~\ref{model:rq1} ($\beta = -0.109$). We additionally
required every specification to meet minimum measurement-quality bars
(formative outer loadings $\geq 0.70$ where applicable, variance inflation
factor $< 3$) and computed the standardized root mean square residual
(SRMR) across all specifications to confirm that model fit did not degrade
as adjustments were added. On this evidence, the negative
compensation--tolerance association is not an artefact of any single
analytical choice. Appendix~\ref{app:vibration} shows an independent
re-computation of the same exercise on observed (mean-score) variables.

\subsection{Multi-group analysis}
\label{robustness:mga}

The second robustness exercise directly compares the structural model
between the 277 residents with compensation experience (CE0 = 1) and the
348 without (CE0 = 0), using permutation-based multi-group analysis
(1{,}000 permutations). A permutation test makes no distributional
assumption and is appropriate for the two groups' unequal sizes; it tests
whether path coefficients are statistically different across groups, going
beyond an earlier, purely visual sub-group comparison. Before comparing any
path coefficients across groups, we first confirmed that discriminant
validity (HTMT) holds \emph{within} each group separately: an HTMT violation
within a subgroup would make that subgroup's path coefficients unreliable,
and any between-group comparison built on them meaningless. Discriminant
validity held in both groups.

We ran two versions of the comparison: MGA~1, the full model with all
background controls; and MGA~2, with sociodemographic controls removed, as a
confounder check -- if a path difference in MGA~1 disappeared once
demographics were added back and a \emph{different} path emerged once they
were removed, that would suggest the demographics were themselves confounding
the group comparison. In fact, the significant findings in MGA~1 and MGA~2 do
not overlap, and no previously-significant path becomes non-significant when
demographics are added or removed; the confounder check passes.

The core finding, robust across both versions, is one of \emph{structural
invariance}: the model's path coefficients do not differ significantly
between compensated and uncompensated residents for almost every path,
including the intangible-cost -- tolerance path, which is significant in both
groups ($p < .001$ in each). One path is the exception: in MGA~2 (without
demographic controls), conflict experience with all species has a
significantly different relationship with tolerance across the two groups
($\beta = -0.187$ for compensated residents versus $\beta = +0.038$ for
uncompensated residents; $\Delta\beta = -0.225$, $p = .039$). We interpret
this as suggestive evidence that compensation experience may amplify the
conflict--tolerance link: among compensated residents, more conflict
experience associates with markedly lower tolerance, whereas among
uncompensated residents the association is close to zero. We supplemented
the permutation $p$-value, which only tells us whether a difference is
statistically detectable, with Cohen's $f^2$ effect size, which asks how
functionally important each predictor is \emph{within} each group
separately (thresholds: negligible $< 0.02$, small $\geq 0.02$, medium
$\geq 0.15$, large $\geq 0.35$; \citealp{Cohen1988}). In MGA~2, intangible
cost's effect size is medium-to-large among uncompensated residents
($f^2 = 0.258$) and small-to-medium among compensated residents
($f^2 = 0.103$) -- the same pattern of intangible cost mattering more for
uncompensated residents holds across both MGA versions and both ways of
quantifying "matters" (statistical difference and effect size).

Finally, the two groups differ in how much of tolerance the model explains
overall: $R^2 = 0.423$ for uncompensated residents versus $R^2 = 0.299$ for
compensated residents in MGA~1. The model captures uncompensated residents'
tolerance more completely than it captures compensated residents' tolerance.
Read alongside the mediation result in Section~\ref{model:rq1} -- where the
measured mediators account for only 18\% of the compensation--tolerance
association -- this lower $R^2$ points the same direction: compensated
residents' tolerance appears to be shaped by additional factors this model
does not measure, plausibly related to their compensation experience itself
(for example, perceptions of payment adequacy, timeliness, or fairness,
which our questionnaire only captured for a subset of applicants and could
not include in the main structural model; Section~\ref{data:quality}) or to
prior conflict history not captured by our conflict-experience items. This
is precisely the kind of open question a longitudinal or better-instrumented
follow-up study, discussed in Section~\ref{discussion}, would be positioned
to resolve.

To summarize RQ3: the model is structurally stable across compensation
groups on nearly every path, and the negative compensation--tolerance
association from RQ1 does not depend on any single analytical choice among
the 8{,}192 tested. The one genuine group difference we find -- conflict
experience mattering more for compensated residents -- and the lower
explanatory power of the model for that same group, both point toward
unmeasured factors specific to the compensation experience as the most
promising direction for future work.
